\scanpage{57}
\deleted{de façon naturelle un revêtement étale
\deleted{\unsure{fini de degré|$T_{\mathrm{sép}}$} \unsure{sép.} de $S$, appelé
$\underline A(T)$ ou} $\overline T=f_\ast^{\mathrm{ét}}(T)$, qui dépend foncto\uline{riellement}
de $T$, \deleted{et transforme} le foncteur $f_\ast^{\mathrm{ét}}$ transforme
d'ailleurs sommes en \deleted{produits} sommes.

On va définir un $S$-morphisme
\[
T\longrightarrow\overline T
\]
i.e. un \unsure{élt|élément} de $\operatorname{Hom}_S(T,\overline T)$, i.e. une
\uline{section} de $\underline{\operatorname{Hom}}_S(T,T$}

Plus \unsure{généralement|précisément}, pour $T$ donné, $S'\leadsto
\underline{\operatorname{Hom}}_S(T,S')$ est un foncteur des variétés
étales de $S$ dans les \unsure{var.|rev.} étales de $S$,
\uline{exact à gauche}. \deleted{Si on choisit un
pt géométrique $a$ de $S$, supposé connexe,}
\inserted{Supposons $S$ connexe.}
\deleted{On obtient ainsi un foncteur
Un théorème général nous
prouve que le foncteur est de la
forme $\varinjlim\operatorname{Hom}($}
% slanted
\margin{(\unsure{inutile?|i.e. la fonctorialité}) \unsure{raisonner|raisonnement|revient à un élément} directement sur le foncteur $S'\leadsto\operatorname{Hom}_S(T,S')$]}
Fixons un
$\mathrm p^{\mathrm t}$ \uline{géométrique} $a$ de $S$ supposé
connexe (ce ne devrait pas
être nécessaire) \deleted{\unsure{Un th. gé|Alors le foncteur}} On obtient
% archive sheet number: 56
